% ============================================================
%  Experiment 6 --- Interfacing FPGA to Digital-to-Analog Converter
%  (renumbered from the former Experiment 10)
% ============================================================
\experimentheader{6}{Experiment 6}{Interfacing FPGA to Digital-to-Analog Converter}{FPGA Lab}

\begin{objectives}
  \begin{enumerate}
    \item To understand the operation of a digital-to-analog converter (DAC).
    \item To interface an FPGA, via the Altera DE2 board, to a DAC.
  \end{enumerate}
\end{objectives}

\begin{equipment}
  \begin{enumerate}
    \item Altera Quartus II software
          \begin{itemize}
            \item version 9/10/12 for Cyclone II
            \item version 9/10/12/latest for Cyclone IV
          \end{itemize}
    \item Altera DE2 board
    \item Multimeter / oscilloscope
    \item DAC0800 digital-to-analog converter
    \item Resistors: $10\,\text{k}\Omega$ (\texttt{$\times$4}), $100\,\Omega$ (\texttt{$\times$1})
    \item Capacitor: $0.01\,\mu\text{F}$
  \end{enumerate}
\end{equipment}

\begin{introduction}
  A digital-to-analog converter (DAC) converts digital pulses into an analog signal. The first criterion for judging a DAC is its resolution, which depends on the number of binary inputs: the number of distinct analog output levels equals $2^n$, where $n$ is the number of data-bit inputs. An 8-bit DAC such as the DAC0800 therefore provides 256 discrete voltage (or current) output levels.

  A DAC can be used to generate a sine wave. The output voltage for a given angle $\theta$ is:
  \[
    V_{out} = 5\,\text{V} + \left(5\,\text{V} \times \sin\theta\right)
  \]
  To find the digital value to send to the DAC for a given angle, multiply the resulting $V_{out}$ by 25.6 (since there are 256 steps across a $10\,\text{V}$ full-scale output).
\end{introduction}

\begin{procedure}

  \textbf{Activity 1 --- DAC0800 operation}\vspace{2pt}

  \textit{Objective: understand the operation of the DAC0800.}

  \begin{enumerate}
    \item Connect the DAC0800 as shown in Figure~\ref{fig:exp10-circuit}.
    \item Apply the digital input values at D0--D7 as listed in Table~\ref{tab:exp10-result}.
    \item Measure the resulting analog output and tabulate your results in Table~\ref{tab:exp10-result}.
  \end{enumerate}

  \begin{boxedfigure}
    \centering
    \includegraphics[width=\linewidth]{figures/Exp6_Figure_1.png}
    \captionof{figure}{DAC0800 test circuit with op-amp output stage.}
    \label{fig:exp10-circuit}
  \end{boxedfigure}

  \needspace{4cm}
  \begin{boxedtable}
    \captionof{table}{DAC0800 digital-to-analog conversion results.}
    \label{tab:exp10-result}
    {\renewcommand{\arraystretch}{1.5}
      \centering
      \begin{tabular}{cc}
        \toprule
        Digital value (decimal) & Analog value \\
        \midrule
        0                       &              \\
        10                      &              \\
        60                      &              \\
        80                      &              \\
        100                     &              \\
        128                     &              \\
        180                     &              \\
        200                     &              \\
        220                     &              \\
        255                     &              \\
        \bottomrule
      \end{tabular}
    }
  \end{boxedtable}
  \newpage

  \textbf{Activity 2 --- Generating a square wave}\vspace{2pt}

  \textit{Objective: interface the FPGA to the DAC0800 and generate a square wave.}

  \begin{enumerate}
    \item Connect the circuit as shown in Figure~\ref{fig:exp10-circuit}.
    \item Connect D0--D7 of the DAC to general-purpose I/O pins on the DE2 board.
    \item Write Verilog code to generate a square wave with a frequency of 100\,Hz and a 70\% duty cycle.
    \item You may use a lookup table, following the same approach as the LCD program in Experiment 5.
    \item Print and submit the Verilog code and simulation result.
  \end{enumerate}

  \textbf{Activity 3 --- Generating a sine wave}\vspace{2pt}

  \textit{Objective: interface the FPGA to the DAC0800 and generate a sine wave.}

  \begin{enumerate}
    \item Using the equation in the Introduction, complete Table~\ref{tab:exp6-sine}.
    \item Connect the circuit as shown in Figure~\ref{fig:exp10-circuit}.
    \item Connect D0--D7 of the DAC to general-purpose I/O pins on the DE2 board.
    \item Write Verilog code to generate a sine wave with a frequency of 100\,Hz.
    \item You may use a lookup table, following the same approach as the LCD program in Experiment 5.
    \item Print and submit the Verilog code and simulation result.
  \end{enumerate}

  \begin{boxedtable}
    \captionof{table}{Angle vs.\ voltage magnitude for the sine wave.}
    \label{tab:exp6-sine}
    {\renewcommand{\arraystretch}{1.45}
      \centering
      \begin{tabular}{cccc}
        \toprule
        Angle ($^\circ$) & $\sin\theta$ & $V_{out}$ (V) & Value sent to DAC ($V_{out}\times 25.6$) \\
        \midrule
        0                &              &               &                                          \\
        30               &              &               &                                          \\
        60               &              &               &                                          \\
        90               &              &               &                                          \\
        120              &              &               &                                          \\
        150              &              &               &                                          \\
        180              &              &               &                                          \\
        210              &              &               &                                          \\
        240              &              &               &                                          \\
        270              &              &               &                                          \\
        300              &              &               &                                          \\
        330              &              &               &                                          \\
        360              &              &               &                                          \\
        \bottomrule
      \end{tabular}
    }
  \end{boxedtable}

\end{procedure}
